Hugging Face - https://huggingface.co/models


# AI Roadmap (After Python)

# Once you're comfortable with Python, move to AI-specific topics in this order:

# NumPy
# Pandas
# Matplotlib
# Scikit-learn
# Linear Algebra basics
# Calculus basics
# Probability & Statistics
# PyTorch
# Neural Networks
# CNN
# RNN / LSTM
# Attention Mechanism
# Transformers
# GPT Architecture
# Train your own small language model



Basics of AI

1.AI vs ML vs DL --------------------------------------------------------------

Artificial Intelligence (AI)
        │
        └── Machine Learning (ML)
                 │
                 └── Deep Learning (DL)

Artificial Intelligence:----------
   -A broad field where machines solve problems intelligently.           
Example:         
   -Chess AI
   -Rule-based chatbots

Machine Learning:-------------
   -The computer learns from data instead of following only hand-written rules.
                 
Example:     
Photos of Cats 🐱
    -↓
    -Train
    -↓
    -Model
    -↓
    -Predict New Cat Images

Deep Learning:----------------
   -A branch of ML that uses neural networks.
Examples:
   -ChatGPT
   -Image generation
   -Speech recognition





2.How AI Learns--------------------------------------------------------------
Imagine teaching a child to recognize apples.

🍎 Apple
🍎 Apple
🍎 Apple
🍎 Apple

After seeing many examples:

🍎

The child says:

"That's an apple."

AI works similarly.


3.Data or Information --------------------------------------------------------------
Text,Images,Number,Audio,Video,Letters ect.. All these are data or information.



4. What is a Dataset?
--------------------------------------------------------------

A dataset is a collection of related data used to train, validate, or test an AI model.

Examples:
- A dataset of cat images.
- A dataset of car images.
- A dataset of numbers.
- A dataset of handwritten letters.
- A dataset of documents.
- A dataset of emails.
- A dataset of SMS messages.
- A dataset of phone call recordings.
- A dataset of customer information.
- A dataset of weather data.

Example Dataset:

| Image      | Label |
|------------|-------|
| cat1.jpg   | Cat   |
| cat2.jpg   | Cat   |
| dog1.jpg   | Dog   |
| dog2.jpg   | Dog   |

Or

| Age | Income | Purchased |
|-----|--------|-----------|
| 25  | 50000  | Yes       |
| 30  | 70000  | No        |
| 22  | 45000  | Yes       |



5.Features and Labels --------------------------------------------------------------
This is one of the most important AI concepts.
Features are the input variables (properties or characteristics) used to make a prediction.

Labels are the correct output (target) that the model learns to predict.

Example 1: Predict House Price
Features:
- Area
- Number of Bedrooms
- Age of House

Label:
- House Price



6.Model Training --------------------------------------------------------------
A model is a mathematical representation of the relationship between features and labels.
Training means teaching the model.
Data
   ↓
Training
   ↓
Model Learns

Example:-

1000 Cat Images
  ↓
Training
  ↓
Model

7.After Training --------------------------------------------------------------
The model is now ready to make predictions.

Example:-

New Cat Image
  ↓
Prediction
  ↓
Predicted Cat Image


8.Model Evaluation --------------------------------------------------------------
After training, the model needs to be evaluated to see how well it performs.

Example:-
Dataset
      ↓
Training
      ↓
Model
      ↓
Prediction






Modules Concepts --------------------------------------------------------------

NumPy
Pandas
Matplotlib
Scikit-learn
Linear Algebra
Calculus
Probability & Statistics
PyTorch
Neural Networks
CNN
RNN / LSTM
Attention Mechanism
Transformers
GPT Architecture
Train Your Own Small Language Model

AI Learning Roadmap - Quick Notes

1.  NumPy

-   Fast numerical computing library.
-   Used for arrays, matrices, and mathematical operations.
-   Think: Super-fast math for Python.

2.  Pandas

-   Library for working with tabular data.
-   Used to read, clean, filter, and analyze CSV/Excel data.
-   Think: Excel inside Python.

3.  Matplotlib

-   Library for creating charts and graphs.
-   Used to visualize data.
-   Think: Graph and chart generator.

4.  Scikit-learn

-   Machine Learning library.
-   Provides ready-to-use ML algorithms such as classification,
    regression, and clustering.
-   Think: Machine Learning toolkit.

5.  Linear Algebra

-   Mathematics of vectors and matrices.
-   Foundation of neural networks.
-   Think: The language of AI mathematics.

6.  Calculus

-   Mathematics of change and optimization.
-   Used to update model weights during training.
-   Think: How AI learns from mistakes.

7.  Probability & Statistics

-   Mathematics of data and uncertainty.
-   Includes mean, median, variance, probability, and distributions.
-   Think: Understanding data.

8.  PyTorch

-   Deep Learning framework.
-   Used to build and train neural networks.
-   Think: Engine for building AI models.

9.  Neural Networks

-   Brain-inspired models that learn patterns from data.
-   Think: The brain of AI.

10. CNN (Convolutional Neural Network)

-   Neural network designed for image processing.
-   Used for image classification and object detection.
-   Think: AI for images.

11. RNN / LSTM

-   Neural networks for sequential data.
-   Used for text, speech, and time-series data.
-   Think: AI with memory.

12. Attention Mechanism

-   Helps the model focus on the most important parts of the input.
-   Think: AI learns what to pay attention to.

13. Transformers

-   Modern neural network architecture based on attention.
-   Used in ChatGPT, translation, summarization, and code generation.
-   Think: Foundation of modern AI.

14. GPT Architecture

-   Generative Pre-trained Transformer.
-   A Transformer model trained to generate human-like text.
-   Think: AI text generator.

15. Train Your Own Small Language Model Steps:

-   Collect Data
-   Clean Data
-   Tokenize
-   Train the Model
-   Save the Model
-   Generate Text

Think: Teaching a computer to understand and generate language.

Roadmap: Python ↓ NumPy ↓ Pandas ↓ Matplotlib ↓ Scikit-learn ↓ Linear
Algebra ↓ Calculus ↓ Probability & Statistics ↓ PyTorch ↓ Neural
Networks ↓ CNN ↓ RNN / LSTM ↓ Attention Mechanism ↓ Transformers ↓ GPT
Architecture ↓ Train Your Own Small Language Model



1. Model Training ————————————————————————————————

Data & Dataset
Tokenizer
Tokens & Embeddings
Nural Network (Model)
Weights & Bias
Forward Pass
Loss Function
Backpropagation
Optimizer
Learning Rate
Batch
Epoch
Training
Validation
Testing
Inference
Fine-tuning
Transformer
Large Language Model (LLM)



1.Data 
 Data is the information used to teach an AI.
 example:- Apple → Fruit , Dog → Animal , Car → Vehicle




2.Dataset
 Dataset is the collection of Data.
 example:- collection of datas like 

 Fruit -> many varieties of fruit
 Animal -> many varieties of animal
 Vehicle -> many varieties of vehicle




3.Tokenizer
 Tokenizer is the function that converts text to numbers.
 example:- 
   like split the text into words and convert them to numbers
   I love AI -> ["I","love","AI"] -> [15,234,98] ai  is not identify the word but it identify the numbers.




4.Tokens & Embeddings
------Tokens:------
Think of it as: A LEGO brick of a sentence.
 A token is a small piece of text produced by the tokenizer.
 example:-
 Sentence:
 I love AI
 
 Tokens:
 ["I", "love", "AI"]


 -----Embeddings:------
 An embedding converts each token ID into a vector (a list of decimal numbers) that captures its meaning.
 example:-
 Text
 "I love AI"
 
         ↓
 Tokenizer
 
 ["I", "love", "AI"]
 
         ↓
 
 [15, 234, 98]        ← Token IDs
 
         ↓
 
 [
  [0.21, -0.45, 0.88, ...],
  [0.73,  0.11, 0.42, ...],
  [-0.19, 0.67, 0.54, ...]
 ]                     ← Embeddings


The AI doesn't work directly with token IDs. It first converts them into embeddings,
because embeddings contain semantic meaning (relationships between words).
Easy Way to Remember
| Concept       | Input    | Output                      |
| ------------- | -------- | --------------------------- |
| **Tokenizer** | Text     | Tokens + Token IDs "I love AI" -> ["I", "love", "AI"] -> [15, 234, 98]          |
| **Token**     | Text     | Small pieces of text "I love AI" + ["I", "love", "AI"]        |
| **Token ID**  | Token    | Integer (e.g., 15, 234, 98) "I" -> 15, "love" -> 234, "AI" -> 98 |
| **Embedding** | Token ID | Vector of decimal numbers [0.21, -0.45, 0.88, ...], [0.73,  0.11, 0.42, ...], [-0.19, 0.67, 0.54, ...] 





5.Neural Network (This is theModel)
Neural Network is the model that learns the relationship between the input and output.

example real time:-
A neural network is the AI's brain.

Data Feeding:-
 user input text is  = "I love AI"
 ↓
 Tokenizer = ["I", "love", "AI"] split the text into words and convert them to numbers
 ↓
 Token IDs = [15, 234, 98] Tokenizer convert the text to token ids
 ↓
 Embeddings = [0.21, -0.45, 0.88, ...], [0.73,  0.11, 0.42, ...], [-0.19, 0.67, 0.54, ...] Embeddings are create the vector of the token ids
 ↓
 Prediction = "I love AI" prediction is the output of the model based on the embeddings

Output:-
This is incorrect because the model's output should be a new prediction, not simply repeat the input.

Input: "I love"
Prediction: "AI"

or

Input: "I love AI"
Prediction: "because it helps people."






6.Layer

A layer is one level inside a neural network. Each layer processes the input and passes it to the next layer.
Think of it as: A factory assembly line.

Example:-
"I love AI"

Input Layer:- 
"I love AI"
 ↓
Layer 1 -Hidden Layer
Understands individual words.
  ↓
Layer 2 -Hidden Layer
Understands the relationship between the words.
  ↓
Layer 3 -hidden layer
Understands the meaning of the sentence.
  ↓
Output Layer:- Positive sentence.

A Layer is a processing stage in a neural network that learns different patterns from the input before passing it to the next layer.






7.Neuron
 A Neuron is a single processing unit inside a layer. It takes the input, processes it, and passes it to the next neuron.
 Think of it as: A single brain cell.

 Example:-
 "I love AI"

 Input Neuron:-
 "I love AI"
 ↓
 Neuron 1
 ↓
 Neuron 2
 ↓
 Neuron 3
 ↓
 Output Neuron:-
 "I love AI"
 ↓
 Output Neuron:-
 "I love AI"




8.Weights 
Weights are the values that are used to make the prediction more accurate.
Inside the neural network there are millions (or billions) of weights 
Weight → Decides how important the input is.

example:-
User Input: "Prakashpathi"
  ↓
Weights:
[0.1, 0.4, 0.7, 1.0]
  ↓
P              → 0.1
Prak           → 0.4
Prakash        → 0.7
Prakashpathi   → 1.0
 ↓
 Prediction: "Prakashpathi"

 "I"    → Less useful for predicting the next word.
 "love" → Gives sentiment (positive feeling).
 "AI"   → Tells what the sentence is about.




9.Bias
Bias → Makes a small adjustment to the final prediction.
Bias is the additional value added to the Weights to make the prediction more accurate.

Example:-
Output = Input × Weight + Bias

User Input:
"I love AI"
        ↓
Neural Network
        ↓
Weights:
"I"    → 0.2
"love" → 0.8
"AI"   → 1.0
        ↓
Bias: 0.3 
+0.3
        ↓
Prediction: "I love AI because it helps people."


-----Relationship Between Forward Pass and Loss Function with Weights and Bias------
10.Forward Pass:-
Forward Pass is the process where the neural network uses the input, weights, and bias to make a prediction.

Example:-
User Input:
"What is 2 + 2?"
        ↓
Neural Network
        ↓
Prediction: "5"




11.Loss Function:-
Loss Function measures how wrong the prediction is by comparing it with the correct answer.

Example:-
Correct Answer:
4
Prediction:
5
        ↓
Loss Function
        ↓
Error = 1
The AI now knows it made a mistake.
---------------------------------------------------------------

13.Backpropagation:-
Backpropagation finds where the AI made a mistake and sends the error backward/Reverse through the network.

Example:- (It tells every layer how to improve.)

Input - Input Layer
↓
Layer 1 - Hidden Layer
↓
Layer 2 - Hidden Layer
↓
Output (Wrong) - Output Layer
↑ - Backward Pass
Error goes backward through the network.

Example after Backpropagation (optimizer updates the weights and bias to reduce the error):-
Wrong Output
↓
Reverse
↓
Layer2
↓
Layer1
↓
Fix Error - Fix the error in the network.

14.Optimizer:-
An Optimizer updates the weights and bias to reduce the error and make future predictions more accurate.
like a Teacher to improve the student's performance.

Example:-
Question:
2 + 2 = ?
        ↓
Prediction:
5
        ↓
Loss Function:
Error = 1
        ↓
Optimizer:
Updates the weights and bias
        ↓
Next Prediction:
4 ✓


Easy way to remember
1.Forward Pass → Makes a prediction.
2.Loss Function → Finds the error.
3.Backpropagation → Finds where the error is.
4.Optimizer → Fixes the weights and bias.
One-line note

Optimizer is an algorithm that updates the model's weights and bias to reduce the prediction error during training.



15. Learning Rate
The learning rate controls how much the AI changes after each mistake.(It controls the speed of the AI learning.)

 Example:
 
 Learning Rate = 0.1 (Fast Learning)
 → Large updates (Learns faster)
 
 Learning Rate = 0.001 (Slow Learning)
 → Small updates (Learns slower)
 
 If the learning rate is too high:
 - The AI may learn too quickly.
 - It may overshoot the correct solution and make inaccurate predictions.
 
 If the learning rate is too low:
 - The AI learns very slowly.
 - It takes more time to reach the correct solution.





16.Epoch:-
Epoch is the number of times the AI sees the entire dataset.
Train full dataset once is one Epoch. Train full dataset 10 times is 10 Epochs.

Example:-
Epoch: 1000 Images
   ↓
After seeing all 1000 images once: Epoch: 1
   ↓
After repeating this 10 times: Epoch: 10

Real-life analogy:-
Imagine you're preparing for an exam.

Book = Dataset

Read the book once  → Epoch 1

Read the book again → Epoch 2

Read the book again → Epoch 3




17.Training:-

Training is the process of teaching the AI the relationship between the input and output.
The AI repeats these steps many times:
Input
↓
Prediction
↓
Calculate Loss
↓
Backpropagation
↓
Update Weights
↓
Repeat

Eventually, the AI becomes better at the task.




18.Validation / Testing:-
Validation is the process of checking how well the AI is performing on the dataset.
like a teacher checking the student's performance. if is learning properly or not.

Example:-
Traini the 100 Image
   ↓
Validation: New 100 Images 
   ↓
Predict the new 100 Images
   ↓
Check the accuracy of the new 100 Images
   ↓
If the accuracy is good, then the AI is learning properly.
   ↓
If the accuracy is not good, then the AI is not learning properly.
   ↓
Repeat the process until the AI is learning properly.



19.Inference:-
 UI for user side.Inference is the process of using the trained model to make a prediction.
 Example:-
 User Input: "I love AI"
   ↓
 Prediction: "I love AI because it helps people."
   ↓
 Output: "I love AI because it helps people."



20.Fine-tuning:-
Definition:-
Fine-tuning is the process of taking a pre-trained AI model and training it on a smaller, specific dataset so it becomes better at a particular task or domain.

Think of it as:
Teach an already smart AI a new skill.



21.Transformer:-
Important Note:-
Transformer is the architecture of the neural network.
The Transformer is not a separate AI model. It is the architecture of the neural network.

Example:-
1.ChatGPT = AI Model
2.Transformer = The architecture used to build ChatGPT
3.A Transformer is a neural network architecture that processes all the words in a sentence together,
learns their relationships, and generates accurate predictions.
-Unlike older models, it doesn't read words strictly one by one. It can consider the whole sentence together.

Real-life Analogy:-
Imagine you're reading a sentence.
  -The cat sat on the mat.

You don't understand it by looking at only one word.
Instead, you understand the whole sentence.
A Transformer works in a similar way.


Neural Network = General category.
Transformer = One specific type of neural network.

Example:-
Neural Network
│
├── CNN (Images) -Convolutional Neural Network
├── RNN (Sequential Data) -Recurrent Neural Network
└── Transformer (Modern LLMs) -Transformer is the architecture of the neural network.




22.Large Language Model (LLM):-

Definition:-

A Large Language Model (LLM) is a Transformer-based AI model trained on a huge amount of text data to understand and generate human language.
1.It can understand and generate human language.
2.All LLM models are Transformer-based.






----------------------------------------------------------------------------------------------------
Complete Flow of Concepts:- (AI Model)

Data
        ↓
Dataset
        ↓
Tokenizer
        ↓
Tokens
        ↓
Token IDs
        ↓
Embeddings
        ↓
Neural Network (Transformer)
        ↓
Weights & Bias
        ↓
Forward Pass
        ↓
Prediction
        ↓
Loss Function
        ↓
Backpropagation
        ↓
Optimizer
        ↓
Update Weights & Bias
        ↓
Learning Rate
        ↓
Repeat (Epochs)
        ↓
Trained Model
        ↓
Fine-tuning (Optional)
        ↓
Large Language Model (LLM)
        ↓
Inference user input
        ↓
User Response = Output





Create your own AI Model:-

=============================
Tokenizer Levels
=============================

Phase 1 - Basic Tokenizers
--------------------------
1. Character-Level Tokenizer
2. Whitespace Tokenizer
3. Word-Level Tokenizer
4. Sentence-Level Tokenizer
5. Regex Tokenizer

Phase 2 - Vocabulary
--------------------
6. Vocabulary
7. Token IDs
8. Encode
9. Decode
10. Special Tokens
    - <PAD>
    - <UNK>
    - <BOS>
    - <EOS>

Phase 3 - Modern Tokenizers
---------------------------
11. Byte-Level Tokenizer
12. Subword-Level Tokenizer
13. Byte Pair Encoding (BPE)
14. WordPiece
15. SentencePiece
16. Unigram Tokenizer

Phase 4 - Build Your Own Tokenizer
----------------------------------
17. Train a Tokenizer
18. Save Vocabulary
19. Load Vocabulary
20. Custom Tokenizer